\chapter{Walking to It: The Dual Active-Set Method}
\label{ch:walking}

This chapter derives the method of Goldfarb and Idnani
\cite{goldfarb1982,goldfarb1983}, which for forty years has been the method of
choice for dense quadratic programs of small to moderate size---$n$ from a handful
to a few thousand. It is often presented as a sequence of updates, which makes it
look more intricate than it is. Read from its invariants it is short: one matrix
is carried, two identities define it, and \emph{everything} the iteration computes
is a consequence of those two lines.

The method's distinguishing property is that it needs no phase-one problem. A
primal active-set method maintains a feasible iterate and works toward
stationarity, so it must begin somewhere feasible, and finding a feasible point is
a linear program in its own right. The dual method reverses the roles: it maintains
stationarity and the multiplier signs throughout, and works toward feasibility.
Its starting point is free, because with no constraint active the stationarity
condition reads $Gx = a$ and the multiplier vector is empty, so
$x = G^{-1}a$ satisfies everything the method asks (Remark~\ref{rem:cold}). One
Cholesky factorisation and one triangular solve replace the phase-one problem
entirely.

The price is that the iterate is infeasible until the last step, so the method
cannot be stopped early for an approximate answer. In exchange it terminates at an
\emph{exactly} feasible point of the active-set lattice rather than approaching one
asymptotically, and it warm-starts almost perfectly---which is why
Chapter~\ref{ch:warmstart} is built on it.

\section{The matrix the solver carries}
\label{sec:factorisation}

Proposition~\ref{prop:sub} gives the subproblem solution in closed form, and one
could implement the method by evaluating it afresh every iteration. That costs a
$k\times k$ factorisation of the dual Hessian per step, $O(nk^2 + k^3)$, and over
the $O(n)$ steps a dual method takes it comes to $O(n^4)$. The whole point of the
Goldfarb--Idnani formulation is to replace this with $O(n^2)$ per step, and it does
so by carrying a single matrix and updating it.

\subsection{Two invariants}

Let $G = R_G\T R_G$ be the Cholesky factorisation with $R_G$ upper triangular, and
set $J_0 := R_G^{-1}$. Then $J_0 J_0\T = G^{-1}$ and $J_0\T G J_0 = I$. The
solver's state is a matrix $J \in \R^{n\times n}$ and an upper triangular
$R \in \R^{k\times k}$ satisfying
\index{invariants (Goldfarb--Idnani)}
\begin{align}
  J\T G J &= I_n, \label{eq:inv1}\\
  J\T A &= \begin{bmatrix} R \\ 0 \end{bmatrix}. \label{eq:inv2}
\end{align}

Read these slowly; everything below is a consequence. Equation~\eqref{eq:inv1} is
equivalent to $JJ\T = G^{-1}$ with $J$ nonsingular, and says the columns of $J$ are
a basis of $\R^n$ orthonormal in the inner product
$\inner{v}{w}_G = v\T G w$. Equation~\eqref{eq:inv2} says that in that basis the
working-set normals are triangular. Partition
\begin{equation}
  J = \begin{bmatrix} J_1 & J_2 \end{bmatrix},
  \qquad J_1 \in \R^{n\times k},\quad J_2 \in \R^{n\times(n-k)};
  \label{eq:split}
\end{equation}
the lower block of~\eqref{eq:inv2} reads $J_2\T A = 0$.

At the cold start $\Active = \emptyset$, so \eqref{eq:inv2} is vacuous and
$J = J_0$ serves. The updates of Section~\ref{sec:updates} maintain both invariants
by post-multiplying $J$ with orthogonal matrices, exactly as
Section~\ref{sec:updating} of Chapter~\ref{ch:freeblock} set up.

\begin{proposition}[What the blocks mean]
\label{prop:blocks}
Suppose \eqref{eq:inv1}--\eqref{eq:inv2} hold with $R$ nonsingular. Then
\begin{enumerate}
  \item[(i)] $R\T R = A\T G^{-1} A$, so $R$ is a Cholesky factor of the dual
        Hessian, unique up to the signs of its rows;
  \item[(ii)] $\range(J_2) = \nullsp(A\T)$, and the columns of $J_2$ are a
        $G$-orthonormal basis of it;
  \item[(iii)] $J_1 = G^{-1} A R^{-1}$ and $\range(J_1) = \range(G^{-1}A)$;
  \item[(iv)] $J_1 J_1\T = G^{-1}A(A\T G^{-1}A)^{-1}A\T G^{-1}$ and consequently
        \begin{equation}
          J_2 J_2\T = G^{-1} - G^{-1}A\bigl(A\T G^{-1}A\bigr)^{-1}A\T G^{-1}.
          \label{eq:reduced}
        \end{equation}
\end{enumerate}
\end{proposition}

\begin{proof}
(i) $A\T G^{-1} A = A\T J J\T A = R\T R$ by~\eqref{eq:inv2}, and two upper
triangular matrices with the same Gram matrix differ by a left factor
$\diag(\pm1)$.

(ii) $J_2\T A = 0$ places $\range(J_2) \subseteq \nullsp(A\T)$, with equality by
dimension since $J$ is nonsingular and $A$ has full column rank; the basis is
$G$-orthonormal because $J_2\T G J_2 = I_{n-k}$ is the trailing block
of~\eqref{eq:inv1}.

(iii) From~\eqref{eq:inv1}, $J^{-1} = J\T G$ and hence $J^{-\top} = G J$.
Inverting~\eqref{eq:inv2}, $A = J^{-\top}[R;0] = G J_1 R$, so
$J_1 = G^{-1}AR^{-1}$, whose range is that of $G^{-1}A$ since $R$ is invertible.

(iv) Substituting $J_1 = G^{-1}AR^{-1}$ and $R\T R = A\T G^{-1}A$,
\[
  J_1 J_1\T = G^{-1}A R^{-1} R^{-\top} A\T G^{-1}
            = G^{-1}A (A\T G^{-1}A)^{-1} A\T G^{-1},
\]
and \eqref{eq:reduced} follows from $J_1J_1\T + J_2J_2\T = JJ\T = G^{-1}$.
\end{proof}

The matrix in~\eqref{eq:reduced} is the \emph{reduced inverse Hessian} of the
working set: it inverts $G$ on the feasible subspace $\nullsp(A\T)$ and annihilates
the complement. So carrying $J$ means carrying that object in factored form, at no
cost beyond the $n\times n$ array---which turns the search directions of the next
section into single matrix--vector products.
\index{reduced inverse Hessian}

\begin{remark}[Why fold the two factorisations together]
An alternative state is $R_G$ together with an explicit orthogonal factor $Q$ of the
QR factorisation of $R_G^{-\top}A$. That holds the same information ($J = J_0Q$) and
stores no more numbers. But every use of $J$ below would become a triangular solve
followed by an orthogonal transformation---two operations where the folded form has
one, and the solve is the shape that does not reach a tuned kernel. The folded form
costs $J$ its triangularity after the first update, which is exactly the trade:
$n^2$ storage and $2n^2$ flops per product, in exchange for the products being a
plain \texttt{gemv} at every iteration whatever the working set does.
\end{remark}

\subsection{The representation is not unique, and need not be}

Invariants~\eqref{eq:inv1}--\eqref{eq:inv2} do not determine $J$ and $R$. By
Proposition~\ref{prop:blocks}, $R$ is pinned only up to a left factor
$S = \diag(\pm1)$, and $J_2$ may be any $G$-orthonormal basis of $\nullsp(A\T)$, so
it is pinned only up to a right factor $V$ with $V\T V = I$. The freedom is exactly
\begin{equation}
  (J_1, J_2, R) \;\longmapsto\; (J_1 S,\; J_2 V,\; S R).
  \label{eq:freedom}
\end{equation}
A Givens chain and a Householder reflection annihilating the same vector produce
factors differing in sign; a differently ordered sequence of insertions and
deletions produces a different $J_2$ outright. The following says none of it
matters, so an implementation may choose its reduction on grounds of speed alone.

\begin{proposition}[Invariance of the iterates]
\label{prop:invariance}
Let $(J,R)$ and $(\hat J,\hat R)$ both satisfy \eqref{eq:inv1}--\eqref{eq:inv2} for
the same $G$ and $A$. Then for every $n_* \in \R^n$ the quantities
\begin{equation}
  d = J\T n_*, \qquad z = J_2 d_2, \qquad r = R^{-1} d_1
  \label{eq:dzr}
\end{equation}
with $d = (d_1,d_2)$ split at $k$ satisfy $\hat z = z$ and $\hat r = r$. If the two
representations differ by~\eqref{eq:freedom} with $V = I$, the equalities hold
exactly in IEEE~754 arithmetic.
\end{proposition}

\begin{proof}
By Proposition~\ref{prop:blocks}(i) and $A\T J = [R\T\ 0]$ we have
$A\T G^{-1} n_* = R\T d_1$, so
\begin{equation}
  r = R^{-1}d_1 = (R\T R)^{-1} A\T G^{-1} n_*
    = \bigl(A\T G^{-1}A\bigr)^{-1} A\T G^{-1} n_* ,
  \label{eq:rclosed}
\end{equation}
which mentions only $A$, $G$, $n_*$. Likewise $z = J_2 J_2\T n_*$, and
by~\eqref{eq:reduced}
\begin{equation}
  z = \Bigl(G^{-1} - G^{-1}A\bigl(A\T G^{-1}A\bigr)^{-1}A\T G^{-1}\Bigr) n_*,
  \label{eq:zclosed}
\end{equation}
again free of the representation. For the last claim, $\hat R = SR$ gives
$\hat d_1 = S d_1$ and $\hat r = (SR)^{-1}Sd_1 = R^{-1}d_1$, in which the only
floating-point operations beyond those forming $r$ are multiplications by $\pm1$;
these are exact.
\end{proof}

\begin{remark}[What the proposition does and does not settle]
It settles that the iteration's \emph{decisions}---which constraint is most
violated, whether the step is full or partial, which constraint is dropped---are
functions of $z$ and $r$ alone, and those are functions of $(G,A,n_*)$ alone. It
does \emph{not} settle that two implementations follow the same trajectory in
finite precision, because the rounding \emph{within} the reductions differs and the
decisions are discontinuous functions of the computed values. That is an empirical
question. What is checkable against a single implementation is the weaker
statement: that the computed $z$ and $r$ agree with the closed
forms~\eqref{eq:zclosed} and~\eqref{eq:rclosed} to machine precision.
\end{remark}

\section{One iteration}
\label{sec:iteration}

The state at the top of an outer iteration is a working set $\Active$ of size $k$,
the factors $(J,R)$, the subproblem solution $x$ and multipliers $u$, and the
running objective $f(x)$. The invariant maintained throughout is
\begin{equation}
  \text{$(x,u)$ solves the subproblem~\eqref{eq:sub}, and}
  \quad u_i \ge 0 \ \text{ for every inequality } i \in \Active,
  \label{eq:invariant}
\end{equation}
which by the discussion after Proposition~\ref{prop:sub} is stationarity,
complementarity, and dual feasibility together. Only primal feasibility is missing,
and the iteration exists to obtain it.

\subsection{Choosing the entering constraint}

Compute the slacks $s_i = c_i\T x - b_i$. If they pass, $x$ is the solution by
Proposition~\ref{prop:sufficient}. Otherwise select the \emph{most violated}
constraint, measured relative to the norm of its own normal:
\begin{equation}
  v_i = \begin{cases} |s_i|, & i \le m_{\mathrm{eq}},\\ -s_i, & i > m_{\mathrm{eq}},\end{cases}
  \qquad
  i_* \in \argmax_i \; \frac{v_i}{\norm{c_i}},
  \qquad\text{terminating when } \max_i \frac{v_i}{\norm{c_i}} \le 0.
  \label{eq:select}
\end{equation}

The scaling is not cosmetic.

\begin{proposition}[Selection is scale invariant]
\label{prop:scale}
Replacing $(c_i,b_i)$ by $(\gamma c_i, \gamma b_i)$ for $\gamma > 0$ leaves the
feasible set, the solution, and the choice made by~\eqref{eq:select} unchanged.
\end{proposition}

\begin{proof}
The rescaled constraint has the same solution set. Its slack scales as
$s_i \mapsto \gamma s_i$ and its normal as $\norm{c_i}\mapsto\gamma\norm{c_i}$, so
the ratio is unchanged, as are the other constraints' ratios.
\end{proof}

Without the division, a constraint written in basis points rather than in units
would be selected first merely for having been written differently, and the
solver's trajectory---though not its answer---would depend on the caller's choice
of units.

\subsection{The two directions}

Write $n_* = c_{i_*}$, $b_* = b_{i_*}$, and form $d$, $z$, $r$ as
in~\eqref{eq:dzr}. The following lemma is the computational core of the method.

\begin{lemma}[Properties of the directions]
\label{lem:directions}
With $d$, $z$, $r$ as in~\eqref{eq:dzr},
\begin{enumerate}
  \item[(i)] $A\T z = 0$: moving along $z$ preserves the working-set equalities;
  \item[(ii)] $n_*\T z = \norm{d_2}^2 = z\T G z \ge 0$, with equality iff $z = 0$;
  \item[(iii)] $G z = n_* - A r$;
  \item[(iv)] $z = 0$ if and only if $n_* = A r$, i.e.\ iff $n_*$ lies in the span
        of the working-set normals, in which case $r$ holds its coordinates there.
\end{enumerate}
\end{lemma}

\begin{proof}
(i) $A\T z = (J_2\T A)\T d_2 = 0$ by~\eqref{eq:inv2}. (ii)
$n_*\T z = (J_2\T n_*)\T d_2 = d_2\T d_2$, and
$z\T G z = d_2\T(J_2\T G J_2)d_2 = d_2\T d_2$ by the trailing block
of~\eqref{eq:inv1}; $J_2$ has full column rank, so $z = 0$ iff $d_2 = 0$. (iii) By
\eqref{eq:reduced}, $Gz = n_* - A(A\T G^{-1}A)^{-1}A\T G^{-1}n_*$, and the second
term is $Ar$ by~\eqref{eq:rclosed}. (iv) If $d_2 = 0$ then
$J\T n_* = (d_1,0) = (J\T A)r$ and $J$ is nonsingular, so $n_* = Ar$; conversely
$n_* = Ar$ gives $Gz = 0$ by (iii).
\end{proof}

Part (ii) identifies $z$ as an ascent direction for the entering constraint's
slack, which grows at rate $n_*\T z > 0$ per unit step. Part (i) says the step
costs nothing in working-set feasibility. Together with the closed forms the
reading is:
\begin{itemize}
  \item $z$ is the $G$-orthogonal projection of the unconstrained direction
        $G^{-1}n_*$ onto the working set's feasible subspace $\nullsp(A\T)$;
  \item $r$ solves a least-squares problem in the $G^{-1}$ metric,
        \begin{equation}
          r = \argmin_{y \in \R^k} \; (Ay - n_*)\T G^{-1} (A y - n_*),
          \label{eq:lsq}
        \end{equation}
        the best approximation of the entering normal by the working-set normals,
        with residual $z = G^{-1}(n_* - Ar)$.
\end{itemize}

\begin{remark}[Why not call a least-squares routine]
Equation~\eqref{eq:lsq} is a genuine least-squares problem and it is tempting to
hand it to a library. Resist: the factor $R$ of its normal equations is
\emph{already held}, so \eqref{eq:dzr} costs one triangular solve, where a library
call would refactorise from scratch at $O(nk^2)$ every iteration. The same applies
to~\eqref{eq:subsol}. \emph{The closed form is the specification, not the
implementation}---a distinction worth carrying out of this book.
\end{remark}

\subsection{The affine path}

Let $u_*$ denote the entering constraint's own multiplier, starting at $0$.

\begin{proposition}[Dual feasibility along the path]
\label{prop:path}
For $t \in \R$ define
\begin{equation}
  x(t) = x + t z, \qquad u(t) = u - t r, \qquad u_*(t) = u_* + t .
  \label{eq:path}
\end{equation}
Then for all $t$: $A\T x(t) = b_\Active$, and stationarity holds for the working set
augmented by the entering constraint,
$G x(t) - a = A\, u(t) + u_*(t)\, n_*$.
\end{proposition}

\begin{proof}
The first claim is Lemma~\ref{lem:directions}(i). For the second, stationarity at
$t = 0$ reads $Gx - a = Au + u_* n_*$, and by Lemma~\ref{lem:directions}(iii),
$Gx(t) - a = Au + u_*n_* + t(n_* - Ar) = A(u-tr) + (u_*+t)n_*$.
\end{proof}

So the \emph{entire path is stationary}, with the entering multiplier rising at
unit rate and the working-set multipliers falling linearly along $-r$.
Complementarity holds for every constraint except the entering one, whose slack is
still nonzero while its multiplier is positive; that single violation is what the
step removes. Two limits on $t$ follow.

\paragraph{The dual limit $t_1$.} Dual feasibility requires $u_i(t) \ge 0$ for every
inequality $i \in \Active$. The binding indices are those with $r_i > 0$, and
\begin{equation}
  t_1 = \min\Bigl\{\, u_i/r_i \;:\; i \in \Active,\ i > m_{\mathrm{eq}},\ r_i > 0 \,\Bigr\},
  \label{eq:t1}
\end{equation}
with $t_1 = +\infty$ when the set is empty. Equalities are excluded: their
multipliers are unrestricted in sign. This is the classical ratio test.

\paragraph{The primal limit $t_2$.} The entering slack is $s_{i_*} + t\,n_*\T z$
with $s_{i_*} < 0$, reaching zero at
\begin{equation}
  t_2 = \frac{|s_{i_*}|}{n_*\T z} = \frac{|s_{i_*}|}{\norm{d_2}^2},
  \label{eq:t2}
\end{equation}
finite and positive whenever $z \ne 0$, and $t_2 = +\infty$ when $z = 0$, in which
case the primal iterate cannot move at all and only the multipliers change.

\paragraph{The step.} Take $t = \min(t_1,t_2)$. If $t_2 \le t_1$ the step is
\emph{full}: the entering constraint joins the working set with multiplier
$u_* + t_2$ and the factorisation is updated. If $t_1 < t_2$ the step is
\emph{partial}: the multiplier of the binding index reaches zero first, that
constraint leaves, the factorisation is downdated, $z$ and $r$ are recomputed, and
the inner loop repeats with the same entering constraint. If both limits are
infinite the problem is infeasible---see Proposition~\ref{prop:farkas}.

\begin{remark}[Equalities violated from above]
\label{rem:reverse}
An equality may be violated in either direction. When $s_{i_*} > 0$ the slack must
be \emph{decreased}, so the step is taken along $-z$: $t$ ranges over the negative
reals, $t_2 = -|s_{i_*}|/n_*\T z$, the ratio test runs against $-r$, and the
entering multiplier accumulates negatively. Nothing else changes, which is why
Proposition~\ref{prop:path} was stated for all $t \in \R$.
\end{remark}

\subsection{The objective, in closed form}

\begin{proposition}[Objective increment]
\label{prop:increment}
Along the path~\eqref{eq:path},
\begin{equation}
  f(x(t)) - f(x) = t\Bigl(\frac{t}{2} + u_*\Bigr)\, n_*\T z ,
  \label{eq:increment}
\end{equation}
strictly positive whenever $z \ne 0$, $t \ne 0$, and $t$, $u_*$ share a sign (in
particular whenever $t > 0$ and $u_* \ge 0$).
\end{proposition}

\begin{proof}
$f(x + tz) = f(x) + t\, z\T(Gx - a) + \tfrac12 t^2\, z\T G z$. Stationarity and
Lemma~\ref{lem:directions}(i) give $z\T(Gx-a) = u_*\,n_*\T z$;
Lemma~\ref{lem:directions}(ii) gives $z\T G z = n_*\T z$. Positivity follows since
$n_*\T z > 0$ and $t(t/2+u_*) > 0$ under the stated sign condition.
\end{proof}

The increment is exact, so the objective is carried as a running total at the cost
of one multiplication per step instead of an $O(n^2)$ re-evaluation, and it is
expressed in quantities the iteration already holds ($n_*\T z$ is needed for $t_2$
anyway). It handles the reversed step of Remark~\ref{rem:reverse} with no special
case: there $t < 0$ and $u_*$ accumulates negatively, so $t/2 + u_* < 0$ and the
product with $t$ is again positive.

\begin{algorithm}[ht]
\caption{The Goldfarb--Idnani dual active-set method}\label{alg:gi}
\begin{algorithmic}[1]
\Require $G \succ 0$, $a$, $C$, $b$, $m_{\mathrm{eq}}$
\State $J \leftarrow \mathrm{chol}(G)^{-1}$; \ $x \leftarrow G^{-1}a$; \
       $f \leftarrow -\tfrac12 a\T x$; \ $\Active \leftarrow \emptyset$; \
       $k \leftarrow 0$; \ $u \leftarrow ()$
\Loop
  \State $s \leftarrow C\T x - b$, with $s_i \leftarrow 0$ for $i \in \Active$;
         choose $i_*$ by~\eqref{eq:select}
  \State \textbf{if} none \textbf{then} \Return $(x,f,u)$
         \Comment{optimal, by Prop.~\ref{prop:sufficient}}
  \State $n_* \leftarrow c_{i_*}$; \ $u_* \leftarrow 0$
  \Loop
    \State $d \leftarrow J\T n_*$; \ $z \leftarrow J_2 d_2$; \
           $r \leftarrow R^{-1} d_1$
    \State $t_1, i_{\mathrm{del}} \leftarrow$ \eqref{eq:t1}; \
           $t_2 \leftarrow$ \eqref{eq:t2};
           \textbf{if} both infinite \textbf{then stop}: infeasible
    \State $t \leftarrow \min(t_1,t_2)$; \ $x \leftarrow x + tz$; \
           $f \leftarrow f + t(t/2 + u_*)\,n_*\T z$; \
           $u \leftarrow u - tr$; \ $u_* \leftarrow u_* + t$
    \If{$t_2 \le t_1$}
      \State $\Active \leftarrow \Active \cup \{i_*\}$, $k \leftarrow k+1$;
             insert; \textbf{break}
    \Else
      \State $\Active \leftarrow \Active \setminus \{i_{\mathrm{del}}\}$,
             $k \leftarrow k-1$; delete
    \EndIf
  \EndLoop
\EndLoop
\end{algorithmic}
\end{algorithm}

\section{Termination}
\label{sec:termination}

Three questions are separate and best kept so: does the inner loop stop, does the
outer loop stop, and what does it mean when the iteration gets stuck.

\begin{proposition}[The inner loop terminates]
\label{prop:inner}
For a fixed entering constraint with $n_* \ne 0$, the inner loop performs at most
$k+1$ passes, where $k$ is the working-set size on entry.
\end{proposition}

\begin{proof}
Every non-exiting pass removes one constraint, so after at most $k$ of them
$\Active = \emptyset$. Then $J_2 = J$ and $z = JJ\T n_* = G^{-1}n_*$, so
$n_*\T z = n_*\T G^{-1}n_* > 0$; thus $z \ne 0$, $t_2$ is finite, and
$t_1 = +\infty$ because the ratio test ranges over an empty set. The step is full
and the loop exits.
\end{proof}

The excluded case $n_* = 0$ is a column of $C$ that is identically zero. Such a
constraint reads $0 \ge b_*$: either vacuous ($b_* \le 0$) or rendering the problem
infeasible. An implementation scores it as infinitely violated when $b_* > 0$,
routing it to the infeasibility verdict rather than to a division by zero.

\begin{proposition}[Strict ascent and finiteness]
\label{prop:ascent}
Every outer iteration ending in a full step with $t_2 \ne 0$ strictly increases the
objective, hence by Proposition~\ref{prop:dual} the dual value. If every full step
has $t_2 > 0$, the method terminates after finitely many outer iterations.
\end{proposition}

\begin{proof}
The iteration's total increment is the sum of~\eqref{eq:increment} over its inner
passes, each positive since the accumulated $u_*$ and each $t$ share a sign
throughout, and the final term nonzero because $t_2 \ne 0$. A working set
determines $(x,u)$ uniquely by Proposition~\ref{prop:sub}, hence determines the
objective; that value strictly increases from one outer iteration to the next, so
no working set recurs. There are finitely many subsets of $\{1,\dots,m\}$.
\end{proof}

\begin{remark}[Degeneracy]
\label{rem:degeneracy}
The hypothesis fails exactly when a full step of length zero occurs, which requires
$s_{i_*} = 0$: more constraints pass through $x$ than the working set holds. A zero
step then adds one without changing $x$, $u$, or the objective, the ascent argument
goes silent, and a cycle is not excluded. This is primal degeneracy, exactly as
Chapter~\ref{ch:geometry} described it; the classical remedies are lexicographic or
least-index rules.
\end{remark}

\subsection{Infeasibility is a certificate, not a failure}

The interesting case is a stuck iteration: $z = 0$, so the primal cannot move, and
$t_1 = \infty$, so no multiplier limits the step. The classical reading is that the
dual is unbounded and therefore the primal infeasible. That argument needs care---by
Proposition~\ref{prop:path} the whole ray is dual feasible, and
by~\eqref{eq:increment} with $z = 0$ the value is constant. The cleanest statement
avoids the dual entirely.

\begin{proposition}[Farkas certificate]
\label{prop:farkas}
\index{Farkas certificate}
Suppose the iteration reaches a state in which the entering inequality $i_*$ is
violated, $s_{i_*} < 0$, and (a) $z = 0$, and (b) $r_i \le 0$ for every inequality
$i \in \Active$. Then $\Omega = \emptyset$.
\end{proposition}

\begin{proof}
By (a) and Lemma~\ref{lem:directions}(iv), $n_* = Ar = \sum_{i\in\Active} r_i c_i$.
Let $\tilde x \in \Omega$. Splitting $\Active$ into equalities $\mathcal{E}$ and
inequalities $\mathcal{I}$,
\[
  n_*\T \tilde x
  = \sum_{i \in \mathcal{E}} r_i\, c_i\T \tilde x
    + \sum_{i \in \mathcal{I}} r_i\, c_i\T \tilde x
  \;\le\; r\T b_\Active ,
\]
since equality members contribute $c_i\T\tilde x = b_i$ exactly and each inequality
member contributes $r_i c_i\T \tilde x \le r_i b_i$ because $c_i\T\tilde x \ge b_i$
and $r_i \le 0$. On the other hand $A\T x = b_\Active$, so
$r\T b_\Active = r\T A\T x = n_*\T x = s_{i_*} + b_* < b_*$. Combining,
$n_*\T\tilde x < b_*$, contradicting $\tilde x \in \Omega$.
\end{proof}

So the solver's infeasibility verdict is a \emph{proof} rather than an exhausted
search, and the proof is a vector the solver computed for its own purposes.
Hypothesis (a) is $t_2 = \infty$ and (b) is $t_1 = \infty$, so the condition an
implementation tests---neither step limit being finite---is literally the hypothesis
of the proposition. This is the second certificate of the book, and it deserves the
same emphasis as the first: an algorithm that returns ``infeasible'' should return
the reason.

\begin{remark}[The remaining hypothesis is not free in floating point]
\label{rem:spurious}
Proposition~\ref{prop:farkas} also needs the entering constraint to be
\emph{genuinely} violated. Exact arithmetic gives that; floating point does not,
and the gap is reachable by the very substitution Section~\ref{sec:updates} argues
for. A Householder reflection and a Givens chain leave the iterate at slightly
different points, so a constraint one implementation leaves just inside its
boundary another can leave the same distance outside. Reaching the stuck state with
such a constraint selected, a solver must neither enforce it, there being nothing
to enforce, nor conclude infeasibility, the conclusion not following; the remedy is
to set it aside for the current iterate and take the next candidate, discarding that
judgement as soon as the iterate moves. The tolerance separating rounding from a
real violation is comparatively easy to set, because infeasibility is macroscopic:
its size is set by the geometry of the constraints, not by the arithmetic.
\end{remark}

\section{Updating the factorisation}
\label{sec:updates}

The two updates are those of Section~\ref{sec:updating} in
Chapter~\ref{ch:freeblock}, applied to $(J,R)$. Insertion appends the entering
normal, costing $O(n^2)$; deletion chases a subdiagonal out of $R$, costing
$O(nk)$. What is new here is a consequence of the step rules.

\begin{proposition}[Full column rank is maintained]
\label{prop:rank}
The insertion is performed only after a full step, which requires $t_2 < \infty$
and hence $z \ne 0$. Then $\alpha \ne 0$ in~\eqref{eq:householder}, so $R'$ is
nonsingular and $A'$ has full column rank $k+1$.
\end{proposition}

\begin{proof}
$|\alpha| = \norm{d_2}$, nonzero iff $z \ne 0$ by Lemma~\ref{lem:directions}(ii).
$R'$ is triangular with diagonal that of $R$ extended by $\alpha$, hence
nonsingular, and ${R'}\T R' = {A'}\T G^{-1}A'$ is then positive definite, forcing
full column rank.
\end{proof}

\emph{The algorithm therefore never has to test for linear dependence among the
working-set normals: the step rules exclude it.} A constraint whose normal lies in
the span of those already held has $z = 0$ by Lemma~\ref{lem:directions}(iv), so its
step is limited by $t_1$ and it can only be added after enough partial steps have
removed the dependence.

This is the structural advantage that Chapter~\ref{ch:guessing} gives up, and it is
worth pausing on. Chapter~\ref{ch:geometry} identified the rank of $C_\Active$ as
the one thing that can go wrong with a general-constraint working-set system, and
observed that it is a property of the \emph{guess}. Here there is no guess: the
walk cannot reach a dependent set. Rank is protected not by a test but by the
geometry of the step rules---an elegance that is easy to overlook and expensive to
lose.

\section{Summary}

One matrix $J$ is carried, defined by $J\T G J = I$ and $J\T A = [R;0]$. Every
quantity the iteration needs is one matrix--vector product or one triangular solve
away, and every quantity it \emph{consumes} has a closed form in $(G,A,n_*)$ alone,
which is why the representation may be chosen for speed and why two implementations
holding different factors compute the same step. The iteration is one affine path
along which stationarity holds identically; two limits cut it short, and which binds
decides whether a constraint is added or dropped. The objective advances by an exact
closed form, positive in both step directions, so the method is a dual ascent and,
absent degeneracy, finite. When the path is blocked in both directions the vector
$r$ is a Farkas certificate.

\begin{notes}
The method is \cite{goldfarb1982,goldfarb1983}; this presentation, organised around
the two invariants, follows \cite{schmelzer2026quadprog}, the mathematics note
accompanying the \texttt{cvx-quadprog} package \cite{cvxquadprog}. The restriction to $G \succ 0$ was
lifted by Boland \cite{boland1997}, who proves finite termination in the
semidefinite case; Forsgren, Gill and Wong \cite{forsgren2016} prove finite
termination for paired primal and dual active-set methods more generally. The modern
successor is a different factorisation rather than a reimplementation: Arnström,
Bemporad and Axehill \cite{arnstrom2022daqp} carry the same dual walk on recursive
$LDL\T$ updates, handle semidefinite $G$ by proximal-point iterations, and ship it as
library-free C for embedded control. Arnström and Axehill
\cite{arnstrom2022certification} compute the \emph{exact} worst-case iteration count
of active-set methods on parametric QP families---stronger than
Proposition~\ref{prop:ascent} and different in kind, being computational and
per-family rather than general and conditional.
\end{notes}

\begin{exercises}

\begin{exercise}
Verify directly from~\eqref{eq:inv1} that $J^{-1} = J\T G$, and use it to give a
one-line proof that $J$ is nonsingular whenever $G$ is.
\end{exercise}

\begin{exercise}
Show that $\Pi := J_2 J_2\T G$ is idempotent, has range $\nullsp(A\T)$, and is
self-adjoint for $\inner{\cdot}{\cdot}_G$. Conclude that $z = \Pi\,G^{-1}n_*$ as
claimed after Lemma~\ref{lem:directions}.
\end{exercise}

\begin{exercise}
Prove that the multiplier of a rescaled constraint scales as
$\lambda_i \mapsto \lambda_i/\gamma$ under $(c_i,b_i)\mapsto(\gamma c_i,\gamma b_i)$,
and deduce that Proposition~\ref{prop:scale} extends to the sign conditions.
\end{exercise}

\begin{exercise}
Work Algorithm~\ref{alg:gi} by hand on $n = 2$, $G = I$, $a = (0,0)\T$, and the two
constraints $x_1 \ge 1$, $x_1 + x_2 \ge 3$. Record $\Active$, $x$, $u$, $t_1$, $t_2$,
and $f$ at every step. Which step is partial?
\end{exercise}

\begin{exercise}
Show that the total objective increase over a whole run equals
$f(x^\star) - f(G^{-1}a)$, and interpret this in terms of
Proposition~\ref{prop:dual}.
\end{exercise}

\begin{exercise}
Construct a two-variable infeasible instance and trace Algorithm~\ref{alg:gi} to the
stuck state. Exhibit the Farkas certificate $r$ explicitly and verify the inequality
in the proof of Proposition~\ref{prop:farkas}.
\end{exercise}

\begin{exercise}
Give an example in which a constraint is added to the working set and later removed.
Why does Proposition~\ref{prop:ascent} not forbid this?
\end{exercise}

\begin{exercise}[Implementation]
Implement Algorithm~\ref{alg:gi}. Verify Proposition~\ref{prop:invariance}
empirically: at each iteration compute $z$ and $r$ both from the carried factors and
from the closed forms~\eqref{eq:zclosed}--\eqref{eq:rclosed}, and report the largest
discrepancy over a run. Then verify Proposition~\ref{prop:increment} by comparing the
running objective against a fresh evaluation of $f(x)$ at every step.
\end{exercise}

\begin{exercise}[Implementation]
Modify your solver to rescale one constraint by $\gamma \in \{10^{-3},\dots,10^{3}\}$
and record the number of additions and removals for each $\gamma$.
Proposition~\ref{prop:scale} claims only that the first \emph{choice} is unchanged;
what do you observe about the whole trajectory?
\end{exercise}

\end{exercises}
